\section{Higher-degree hypersurface tangent cones}\label{higher-surfaces}
We now blow up $I=\m$ and suppose that
$G=\operatorname{gr}_{\m}R=P/(f_e)$, with $P=k[X_0,\ldots,X_3]$,
$f_e$ homogeneous of degree $e\ge3$, so $R$ is Cohen--Macaulay of dimension three by the preceding section.
With $X,D,L,V,M,Q$ as in Section~\ref{blowup}, put
$E=V|_D$, $W=\Sym^3 E$, and let $b$ denote the number of columns of $U$.
As before, the linear part of $U$ must have full column rank on $D$.
By the trace argument, it suffices to work in residue characteristic three.
On a degree-$e$ surface, $H^2(D,\mathcal O_D(-1))^\vee=G_{e-3}$.
This is the group that vanished in the proof of Theorem C. The following lifting bound and presentation calculations account for it.
\begin{lemma}\label{high-order-cubic-lifting}
Every section of $L^{e+b-2}Q$ has a global preimage in
$L^{e+b-2}M$.
\end{lemma}
\begin{proof}
The hypersurface sequence on $\mathbf P^3_k$ and adjunction give
\[
 H^1(D,\mathcal O_D(n))=0\quad(n\in\mathbf Z),\qquad
 \omega_D=\mathcal O_D(e-4),
\]
and $H^2(D,\mathcal O_D(n))=0$ for $n\ge e-3$.
Relative Serre vanishing for $L$, followed by descending induction in
$0\to L^{n+1}\to L^n\to\mathcal O_D(n)\to0$, gives
\[
 H^1(X,L^n)=0\ (n\ge0),\qquad H^2(X,L^n)=0\ (n\ge e-3).
\]
Twisting \eqref{bundlec} by $L^{e-2}$ now gives
$H^1(X,L^{e-2}V)=0$.
By $V^\vee\simeq V\otimes L^{-b}$, this is the first cohomology of
the kernel in \eqref{cubicsequence} twisted by $L^{e+b-2}$.
The required surjectivity follows from its cohomology sequence.
\end{proof}

\begin{corollary}\label{three-entries}
Suppose $R\subseteq S$ is a module-finite extension with $S/R\simeq\operatorname{coker}U$ and the ideal $I_1(U)$ of entries has three generators. If the linear matrix has full column rank on $D$, then $R\subseteq S$ splits for every $b$ and every $e\ge3$.
\end{corollary}
\begin{proof}
The initial entries span at least three dimensions in $P_1$. Otherwise their common zero locus would contain a line, which meets $D$ over an algebraic closure. Since the entry ideal has three generators, the span has dimension three. Choose generators $a_0,a_1,a_2$ whose linear parts form a basis, and complete these forms to coordinates $X_0,X_1,X_2,Z$ of $P_1$. The point $[0:0:0:1]$ lies outside $D$. Consequently the coefficient of $Z^e$ in $f_e$ is nonzero. It follows that the linear parts form a system of parameters in $G$, and $a_0,a_1,a_2$ form a regular sequence in the Cohen--Macaulay ring $R$. Indeed the associated graded ring of their quotient is a quotient of $G/(X_0,X_1,X_2)$ and has finite length.
Projection from that point gives a finite surjection $D\to\PP^2_k$, as $f_e$ is monic in $Z$ after scaling. The linear matrix on $D$ is the pullback of $\overline U_{I_1(U)}$ on this plane. Its full rank descends to the plane, and Theorem C(i) applies.
\end{proof}

\subsection{Two relations}
\begin{theorem}\label{hyp-two}
Let $b=2$ and suppose $S/R\simeq\operatorname{coker}U$. If $e=3$ or the entries of the $4\times2$ matrix $U$ generate $\m$, then the module-finite extension $R\subseteq S$ splits. The same holds for any perturbation in $\m^2$.
\end{theorem}
\begin{proof}
Write $u,v$ for the columns of the linear part of $U$.
Define the graded $P$-module
\[
 N=\operatorname{coker}\bigl(
 P(-1)^{20}\longrightarrow P^4,\quad
 (A,B)\longmapsto Au+Bv\bigr),
\]
where $A$ and $B$ range over symmetric $4\times4$ matrices, and put
$N_G=N\otimes_PG$.
We first prove
\begin{equation}\label{two-relation-vanishing}
 (N_G)_q=0\qquad(q\ge3)
\end{equation}
under the stated hypotheses.
The standard generators $\epsilon_j$ of $N$ satisfy
$u_i\epsilon_j=-u_j\epsilon_i$ and $u_i\epsilon_i=0$.
Consequently
\[
 u_iu_j\epsilon_l=-u_i u_l\epsilon_j=u_lu_j\epsilon_i=-u_ju_i\epsilon_l.
\]
Since $2$ is a unit, this gives
$I(u)^2N=I(v)^2N=0$.
If the entries of $U$ generate $\m$, their initial forms span $P_1$, so $P_+=I(u)+I(v)$ and $(P_+)^3N=0$.

It remains to treat $e=3$ when those initial forms do not span $P_1$.
Their span then has dimension exactly three. A span of dimension at
most two would give a projective common zero locus of positive
dimension meeting the hypersurface $D$, contradicting the rank
assumption.
Choose coordinates so that all entries lie in
$P_0=k[X_0,X_1,X_2]$ and put $Z=X_3$.
Their common zero is outside $D$, so the coefficient of $Z^e$ in
$f_e$ is nonzero.  Projection from that point gives a finite surjective
map $D\to\mathbf P^2_k$.
The two columns are therefore independent at every geometric point of
$\mathbf P^2$.

Let $N_0$ be the corresponding module over $P_0$.
We show that $(N_0)_1=0$.
Write the columns as linear maps $A_1,A_2:k^3\to k^4$.
By duality against symmetric matrices, a functional on $(N_0)_1$
is represented by a $3\times4$ matrix $B$ for which both $A_iB$
are alternating.
The map $A_1$ is injective, so choose bases with
\[
 A_1=\binom{I_3}{0},\qquad A_2=\binom{H}{w}.
\]
Alternation of $A_1B$ forces $B=(B_0\ 0)$, where $B_0$ is an
alternating $3\times3$ matrix.
Alternation of $A_2B$ then gives $HB_0$ alternating and $wB_0=0$.
If $B_0\ne0$, its rank is two.
Both $B_0$ and $HB_0$ kill $\ker B_0$. The space of alternating
forms on $k^3$ with this property is one-dimensional.
Hence $HB_0=\lambda B_0$.
For any nonzero $v\in\operatorname{im}B_0$, we obtain
$Hv=\lambda v$ and $wv=0$, so $A_2v=\lambda A_1v$, a contradiction.
Thus $B=0$ and $(N_0)_1=0$.

Since $N_0$ is generated in degree zero and has no relations there,
$N_0=k^4$ in that degree and is zero in higher degrees.
It follows that
$N_G\simeq(k[Z]/(Z^e))^4$, after rescaling the coefficient of $Z^e$.
For $e=3$, this proves \eqref{two-relation-vanishing}.

We next translate this graded calculation into a lifting statement.
The degree-three part of the Koszul resolution of the symmetric
quotient on $D$ is
\[
 0\to\mathcal O_D(-2)^4\to\mathcal O_D(-1)^{20}
 \to\mathcal O_D^{20}\to W\to0 .
\]
It is exact because the bundle presentation is locally split, so its
Koszul generators become two independent polynomial variables.
Set $B'=\ker(\mathcal O_D^{20}\to W)$.
Serre duality and the vanishing of $H^1(\mathcal O_D(n))$ give
\begin{equation}\label{two-relation-obstruction}
 H^1(D,B'(t))^\vee=(N_G)_{e-2-t}.
\end{equation}
Indeed, dualizing the map on $H^2$ of the two leftmost terms gives
$(A,B)\mapsto Au+Bv$, with coefficients in degree $e-3-t$.
For $t<0$, the sequence
$0\to B'(t)\to\mathcal O_D(t)^{20}\to W(t)\to0$ identifies
$H^0(D,W(t))$ with $H^1(D,B'(t))$.
For $t\ge0$, it identifies the latter group with the cokernel of
the map from degree-$t$ coefficients times the twenty cubic monomials.
Here $H^0(D,\mathcal O_D(t))=G_t$.

Since $L^nQ|_D=W(n-4)$, for $0\le n<4$ the dual of its section
space is $(N_G)_{e+2-n}$.
Successive restriction therefore puts $\Phi$ in $L^3Q$ when $e=3$,
and in $L^4Q$ when $e\ge4$.
For $e=3,4$, Lemma~\ref{high-order-cubic-lifting} already applies.
For $e\ge5$ and $0\le t\le e-5$, equations
\eqref{two-relation-vanishing} and \eqref{two-relation-obstruction}
show that $H^0(D,W(t))$ is generated by those monomials.
Given a remainder in $L^{4+t}Q$, express its restriction in that form,
lift the coefficients from $G_t=\m^t/\m^{t+1}$ to $\m^t$, and use the
global generators of $V$.
Symmetric multiplication gives a global preimage under $j$ of the
resulting section.  Subtraction raises the remainder to $L^{5+t}Q$.
After finitely many steps it lies in $L^eQ$ and lifts by
Lemma~\ref{high-order-cubic-lifting}.
Lemma~\ref{cubic-lifting} then gives the retraction.
\end{proof}

For an explicit example, take
$R=\mathbf Z_3[[x,y,z]]/(3^e+x^e+y^e+z^e)$, where
$e=4r$, $r\ge1$, and $3\nmid r$.
Use the columns $(3,x,y,z)^t$ and $(x,y,z,3)^t$.
Their linear forms span $P_1$, and their rank drops only at
$[1:\lambda:\lambda^2:\lambda^3]$ with $\lambda^4=1$.
At every such point, the initial Fermat equation has value $4=1$ in
the residue field.
Thus the criterion applies to every module-finite extension with this
presentation, and also to every presentation obtained by perturbing
its entries in order at least two.

\subsection{Three relations and the exceptional section}
We may use a hypersurface equation in a regular ambient ring in this proof. Completion preserves the evaluation ideal by flat base change, and preserves $G$ and the linear matrix. By the Cohen structure theorem, $\widehat R$ is a quotient of a regular local ring \cite{Stacks}. In the kernel choose elements whose linear parts form a basis of the kernel on cotangent spaces. They are part of a regular system of parameters, so quotienting by them gives a regular local ring $(A,\mathfrak n,k)$ of dimension four mapping onto $\widehat R$. Since $\widehat R$ has depth three, Auslander--Buchsbaum gives $\pd_A\widehat R=1$. The nonzero kernel of $A\to\widehat R$ is consequently free of rank one, say $(F)$. The initial ideal of $(F)$ is generated by $\operatorname{in}F$ in $\operatorname{gr}_{\mathfrak n} A=P$. Under the chosen tangent coordinates $\operatorname{in}F$ is a scalar multiple of $f_e$. We can therefore work with $R=A/(F)$ without imposing a presentation of the original local ring.

Write $u,v,w$ for the linear columns of the $5\times3$ matrix $U$.
For symmetric $5\times5$ matrices, define
\begin{align*}
 d(B_1,B_2,B_3)&=(B_2u-B_1v,B_3u-B_1w,B_3v-B_2w),\\
 \delta(\alpha,\beta,\gamma)&=w^t\alpha-v^t\beta+u^t\gamma .
\end{align*}
\begin{theorem}[Three relations]\label{three-relation-criterion}
Suppose the $5\times3$ linear matrix has rank three everywhere on $D$ and its entries span $P_1$. Write
$d_q:(P_{q-1})^{45}\to(P_q)^{15}$ for $d$ in degree $q$. If
\[
                   \operatorname{rank} d_4=469,
\]
every module-finite extension $R\subseteq S$ with $S/R\simeq\operatorname{coker}U$ splits. The smaller test $\operatorname{rank}d_2=130$ suffices as well. For each $D$ the admissible matrices form a geometrically nonempty dense open, with dense rational points when $k$ is infinite.
\end{theorem}
\begin{proof}
Consider the graded complex
\[
 P(-1)^{45}\xrightarrow{d}P^{15}\xrightarrow{\delta}P(1)
\]
with $P^{15}$ in cohomological degree zero. Symmetry gives $\delta d=0$. Write $N_P$ for its middle cohomology and $N_G$ for the middle cohomology after reduction modulo $f_e$. By a constant change of basis, $\delta$ is the row of four variables followed by eleven zeros. Its cokernel is $k(1)$ and its kernel is generated in degrees zero and one by the extra basis vectors and the Koszul relations. Thus $N_P$ is generated in those degrees, and
\[
 \dim_k(\ker\delta)_q=15\dim_k P_q-\dim_k P_{q+1}\quad(q\ge0).
\]
The dimensions are $130$ in degree two and $469$ in degree four. Surjectivity onto the kernel in degree $q\ge2$ forces $(N_P)_j=0$ for every $j\ge q$. In particular the degree-four test follows from the degree-two test. The exact sequence of complexes given by multiplication by $f_e$ yields
\begin{equation}\label{three-relation-obstruction}
 0\longrightarrow N_P/f_eN_P\longrightarrow N_G\longrightarrow k(1-e)\longrightarrow0 .
\end{equation}
Here $f_e$ kills $k(1)$. In degrees $q\ge4$ the only possible nonzero part of $N_G$ is one copy of $k$ when $q=e-1$.

\emph{The exceptional section.}
The degree-three symmetric quotient now has the exact resolution
\[
 0\to\mathcal O_D(-3)\to\mathcal O_D(-2)^{15}\to
 \mathcal O_D(-1)^{45}\to\mathcal O_D^{35}\to W\to0.
\]
Exactness again follows by locally completing the three independent
linear forms to polynomial coordinates.
Set $B'=\ker(\mathcal O_D^{35}\to W)$.
Break the resolution at the image of its second arrow.
The vanishing of $H^1(\mathcal O_D(t))$ and Serre duality on the two
resulting short exact sequences identify
\[
 H^1(D,B'(t))^\vee=(N_G)_{e-2-t}.
\]
the transposed maps on the $H^2$ groups are exactly $d$ and $\delta$.
As before, for $t<0$ this also computes $H^0(D,W(t))^\vee$.

Now $L^nQ|_D=W(n-6)$. For $0\le n<\min(5,e+1)$ its sections are dual to $(N_G)_{e+4-n}$. This index is at least four and greater than $e-1$, so the section space vanishes. Successive restriction puts $\Phi$ in $L^{\min(5,e+1)}Q$. For $e=3,4$, Lemma~\ref{high-order-cubic-lifting} and then Lemma~\ref{cubic-lifting} give the retraction. Assume $e\ge5$. The next restriction lies in $H^0(D,W(-1))$, a one-dimensional space by \eqref{three-relation-obstruction} in degree $e-1\ge4$.

\emph{Lifting the exceptional section.}
Let $Y$ be the blow-up of $\operatorname{Spec}A$ at its closed point,
with exceptional divisor $\Delta=\mathbf P^3_k$.
The surjection of Rees algebras embeds $X$ in $Y$, with
$\mathcal O_Y(-\Delta)|_X=L$ and scheme-theoretic intersection
$X\cap\Delta=D$.
Lift the entries of $U$ to $A$ and divide by the local equation of
$\Delta$.  This gives
$B_1'=\mathcal O_Y(\Delta)^3\to B_0'=\mathcal O_Y^5$.
Its restriction to $X$ is the subbundle in \eqref{bundlec}.

Let $K_i$ be the degree-$i$ part of the Koszul complex of the
corresponding linear forms in the symmetric algebra of $B_0'$.
Place $\operatorname{Sym}^iB_0'$ in degree zero, so that its term in
degree $-j$ is
$\bigwedge^jB_1'\otimes\operatorname{Sym}^{i-j}B_0'$.
Koszul multiplication gives
\[
 (K_1\otimes K_2)(\Delta)\longrightarrow K_3(\Delta).
\]
In degree $-3$, this is
\[
 B_1'\otimes\bigwedge\nolimits^2B_1'
 \longrightarrow\bigwedge\nolimits^3B_1',
\]
twisted by $\Delta$.
It is exterior multiplication from nine copies of
$\mathcal O_Y(4\Delta)$ to one.
A basis tensor maps to a generator with coefficient $1$, so the
map is surjective in characteristic three as well.

For $0\le j\le3$, the line bundle $\mathcal O_Y(j\Delta)$ has no
higher cohomology.
For $j=4$, its only positive cohomology is $H^3$, and restriction is
an isomorphism
\[
 H^3(Y,\mathcal O_Y(4\Delta))
   \simeq H^3(\Delta,\mathcal O_\Delta(-4))=k.
\]
These statements follow first for nonpositive multiples by relative
Serre vanishing and descending induction, and then for $j=1,\ldots,4$
from restriction to $\Delta$ and the cohomology of projective space.

Consider the hypercohomology spectral sequence for
$(K_1\otimes K_2)(\Delta)$ on $Y$.
Its only positive-cohomology term is $k^9$ in bidegree $(-3,3)$.
There is no incoming differential, and every possible outgoing
target either vanishes or lies outside the complex.
The edge map therefore surjects from degree-zero hypercohomology
onto $k^9$.
After restriction to $\Delta$, all terms are sums of
$\mathcal O_\Delta(-j)$ for $1\le j\le4$.
Consequently this $k^9$ is the entire degree-zero hypercohomology
on $\Delta$.
Naturality of the edge maps and the preceding restriction
isomorphism show that restriction from $Y$ onto this group is
surjective.
Multiplication maps it onto the corresponding one-dimensional
hypercohomology of $K_3|_\Delta(-1)$.

To identify the image on $D$, use the short exact sequence of
complexes
\[
 0\to K_3|_\Delta(-e-1)
   \xrightarrow{\ f_e\ }K_3|_\Delta(-1)
   \to K_3|_D(-1)\to0.
\]
The first complex has zero degree-zero hypercohomology.
Indeed, all its twists are negative, so only the top cohomology
of the leftmost term can contribute in total degree zero.
The kernel of the corresponding differential is dual to the
cokernel of
\[
 (P_{e-1})^{15}\xrightarrow{\ \delta\ }P_e,
\]
which is zero because the entries of $\delta$ span $P_1$.
Restriction therefore injects the one-dimensional group
$\mathbb H^0(\Delta,K_3|_\Delta(-1))$ into
$H^0(D,W(-1))$.
The dimension calculation above makes this map an isomorphism.

On $X$, the locally split sequence \eqref{bundlec} makes $K_i|_X$
a locally free resolution of $\operatorname{Sym}^iV$.
Thus the restrictions of the source and target of Koszul
multiplication resolve $L^5M$ and $L^5Q$, respectively.
Naturality along $D\subset X\subset Y$ and $D\subset\Delta\subset Y$
now gives the surjection
\begin{equation}\label{exceptional-cubic-lift}
 H^0(X,L^5M)\xrightarrow{\ j|_D\ }H^0(D,W(-1)).
\end{equation}

\emph{Completion of the lifting.}
Choose $\mu\in H^0(X,L^5M)$ whose image under
\eqref{exceptional-cubic-lift} is the restriction of $\Phi$.
The remainder $\Phi-j(\mu)$ lies in $L^6Q$.
For $0\le t\le e-6$, equation \eqref{three-relation-obstruction}
gives $(N_G)_{e-2-t}=0$, since $4\le e-2-t\le e-2$.
The monomial-lifting argument used for two relations therefore
raises the order to $e+1$.
When $e=5$, order six already suffices.
Lemma~\ref{high-order-cubic-lifting}, with $b=3$, lifts the final
remainder, and Lemma~\ref{cubic-lifting} gives the retraction.
\end{proof}

\begin{proof}[Geometric nonemptiness in Theorem~\ref{three-relation-criterion}]
In the affine space of $5\times3$ matrices of linear forms,
the condition that the entries span $P_1$ is open.
Full rank on $D$ is open by properness of the rank-drop incidence projection.
On the spanning locus, rank $469$ for $d_4$ is open by the minor criterion, and is the maximal possible rank since $\delta d=0$.
Let $\mathcal U_D$ be the intersection of these three opens.

To prove geometric nonemptiness, consider
\[
 U_0=\begin{pmatrix}a&0&0\\ b&a&0\\ c&b&a\\ d&c&b\\0&d&c\end{pmatrix}
\]
in four parameter coordinates.
The unique point where its rank drops is $[0:0:0:1]$. If the first nonzero coordinate
is among the first three, the corresponding three consecutive rows
give a nonzero triangular minor.
Its entries span $P_1$, and Appendix~\ref{cubcert} certifies
$\operatorname{rank}d_2=130$ over $\mathbf F_3$, hence over every
extension of that field.
Over $\bar k$, a linear coordinate change moves its unique bad point
outside the fixed hypersurface $D$.
Such a change preserves all graded ranks. The generation of $N_P$ in degrees zero and one then gives the degree-four test.
Thus $\mathcal U_D$ is geometrically nonempty and dense.

If $k$ is infinite, $\mathcal U_D(k)$ is dense by density of rational
points in affine space over an infinite field.
Every lift of a matrix in $\mathcal U_D(k)$ satisfies
the splitting criterion just proved.
All three conditions depend only on the linear part, so the same is
true after arbitrary perturbations of the entries in $\m^2$.
\end{proof}

For a family over the finite residue field, take
$A=\mathbf Z_3[[x,y,z]]$ with parameters $(3,x,y,z)$ and
$F=3^e+x^e+y^e+z^e$, for any $e\ge3$.
The initial equation takes value $1$ at the bad point of $U_0$, so
that matrix itself lies in $\mathcal U_D(\mathbf F_3)$.
If $3\nmid e$, the Fermat surface $D$ is smooth.
At a nonmaximal prime of $R$, at least one of $x,y,z$ is invertible:
otherwise $F=0$ would force $3$ into the prime as well.
The corresponding partial derivative of $F$ is a unit, proving that
$R$ is regular away from its closed point.
Since $R$ is Cohen--Macaulay, Serre's criterion makes it normal.
These isolated singularities have unbounded multiplicity, and
\[
 h^2(D,\mathcal O_D)=\binom{e-1}{3}.
\]
Thus the criterion does not rely on the vanishing of this exceptional
cohomology.  The exceptional section lifts by \eqref{exceptional-cubic-lift}.

For a general linear presentation on a higher-degree hypersurface, full column rank alone leaves the cubic lifting problem open in this argument.
